\documentclass[
    language=english,
    doctype=article,
    institution=none,
]{../../omnilatex}

\RequirePackage{config/document-settings}
\addbibresource{../../bib/bibliography.bib}

\begin{document}

\maketitle

\begin{abstract}
This document specifies the SecureLink transport protocol, a lightweight,
cryptographically authenticated protocol for secure machine-to-machine
communication in distributed systems.  SecureLink provides mutual
authentication, forward secrecy, and resistance to replay attacks while
imposing minimal computational overhead on constrained devices.
\end{abstract}

\tableofcontents

\section{Scope}

This specification defines the wire format, handshake procedure, session
management, and error handling for SecureLink v1.0.  It is intended for
implementation by:

\begin{itemize}
    \item IoT gateway manufacturers requiring authenticated device-to-cloud
          communication.
    \item Microservice platforms seeking an alternative to mutual TLS with
          lower connection establishment latency.
    \item Industrial control systems where deterministic handshake timing is
          critical.
\end{itemize}

This specification does not cover application-layer semantics, key management
infrastructure, or compliance certification procedures.

\section{Requirements}

The key words \textbf{MUST}, \textbf{SHOULD}, and \textbf{MAY} in this
document are to be interpreted as described in RFC 2119
\cite{rfc2119}.

\subsection{Functional Requirements}

\begin{table}[htbp]
    \centering
    \caption{Functional requirements.}
    \label{tab:func-req}
    \begin{tabular}{@{}llp{5.5cm}@{}}
        \toprule
        \textbf{ID} & \textbf{Priority} & \textbf{Requirement} \\
        \midrule
        FR-01 & MUST & The protocol shall establish a mutually authenticated
            session within two round trips. \\
        FR-02 & MUST & The protocol shall support ephemeral key exchange
            providing forward secrecy. \\
        FR-03 & MUST & The protocol shall detect and reject replayed
            handshake messages. \\
        FR-04 & SHOULD & The protocol shall support session resumption to
            reduce handshake overhead. \\
        FR-05 & MAY   & The protocol shall support multiplexed streams over
            a single connection. \\
        FR-06 & MUST & The protocol shall negotiate cipher suites using a
            server-preference ordering. \\
        \bottomrule
    \end{tabular}
\end{table}

\subsection{Non-Functional Requirements}

\begin{table}[htbp]
    \centering
    \caption{Non-functional requirements.}
    \label{tab:nfr}
    \begin{tabular}{@{}lp{7cm}@{}}
        \toprule
        \textbf{ID} & \textbf{Requirement} \\
        \midrule
        NFR-01 & Handshake completion shall not exceed 10\,ms on an ARM Cortex-M4
            at 80\,MHz with hardware AES acceleration. \\
        NFR-02 & The protocol shall support connections lasting at least 24 hours
            without mandatory re-authentication. \\
        NFR-03 & Bandwidth overhead for the handshake shall not exceed 512 bytes. \\
        NFR-04 & The protocol specification shall be implementable without
            dependency on any single cryptographic library. \\
        \bottomrule
    \end{tabular}
\end{table}

\section{Architecture}

\subsection{Protocol Layers}

SecureLink is structured in three layers:

\begin{enumerate}
    \item \textbf{Transport Layer:} Provides reliable, ordered byte delivery
          over TCP or a reliable UDP variant (QUIC).  The transport layer is
          responsible for fragmentation, reassembly, and congestion control.
    \item \textbf{Security Layer:} Performs the cryptographic handshake,
          derives session keys, encrypts and decrypts application data, and
          manages session lifecycle.
    \item \textbf{Application Layer:} Carries multiplexed logical streams,
          each identified by a 16-bit stream identifier.
\end{enumerate}

\subsection{Handshake Sequence}

The SecureLink handshake proceeds in four phases:

\begin{enumerate}
    \item \textbf{ClientHello:} The client sends a random nonce (32 bytes),
          a list of supported cipher suites, and an optional session ticket
          for resumption.
    \item \textbf{ServerHello:} The server responds with its selected cipher
          suite, a random nonce (32 bytes), its certificate chain, and an
          ephemeral public key.
    \item \textbf{Key Confirmation:} Both parties derive the session keys
          using X25519 key exchange and HKDF.  The client sends a
          \texttt{Finished} message keyed with the derived handshake key.
    \item \textbf{Session Established:} The server verifies the
          \texttt{Finished} message and responds with its own
          \texttt{Finished} message.  Application data may follow immediately.
\end{enumerate}

\section{Interfaces}

\subsection{Wire Format}

All multi-byte integers are encoded in network byte order (big-endian).

\begin{verbatim}
 0                   1                   2                   3
 0 1 2 3 4 5 6 7 8 9 0 1 2 3 4 5 6 7 8 9 0 1 2 3 4 5 6 7 8 9 0 1
+-+-+-+-+-+-+-+-+-+-+-+-+-+-+-+-+-+-+-+-+-+-+-+-+-+-+-+-+-+-+-+-+
|  Version (8)  |  MsgType (8)  |         Length (16)           |
+-+-+-+-+-+-+-+-+-+-+-+-+-+-+-+-+-+-+-+-+-+-+-+-+-+-+-+-+-+-+-+-+
|                        Payload (variable)                     |
+-+-+-+-+-+-+-+-+-+-+-+-+-+-+-+-+-+-+-+-+-+-+-+-+-+-+-+-+-+-+-+-+
|                        MAC (32)                               |
+-+-+-+-+-+-+-+-+-+-+-+-+-+-+-+-+-+-+-+-+-+-+-+-+-+-+-+-+-+-+-+-+
\end{verbatim}

\begin{table}[htbp]
    \centering
    \caption{Message types.}
    \label{tab:msgtypes}
    \begin{tabular}{@{}rl@{}}
        \toprule
        \textbf{Code} & \textbf{Message Type} \\
        \midrule
        0x01 & ClientHello \\
        0x02 & ServerHello \\
        0x03 & Certificate \\
        0x04 & KeyExchange \\
        0x05 & Finished \\
        0x06 & ApplicationData \\
        0x07 & CloseNotify \\
        0xFF & Alert \\
        \bottomrule
    \end{tabular}
\end{table}

\subsection{Cipher Suites}

\begin{table}[htbp]
    \centering
    \caption{Supported cipher suites.}
    \label{tab:ciphers}
    \begin{tabular}{@{}ll@{}}
        \toprule
        \textbf{ID} & \textbf{Algorithm} \\
        \midrule
        0x01 & X25519 + AES-256-GCM + SHA-384 \\
        0x02 & X25519 + ChaCha20-Poly1305 + SHA-256 \\
        0x03 & P-256 + AES-256-GCM + SHA-384 \\
        \bottomrule
    \end{tabular}
\end{table}

\section{Constraints}

\begin{itemize}
    \item \textbf{Maximum record size:} 16\,384 bytes (excluding MAC).
    \item \textbf{Maximum certificate chain depth:} 3.
    \item \textbf{Session ticket lifetime:} 24 hours maximum.
    \item \textbf{Replay window:} 64 messages or 30 seconds, whichever is
          larger.
    \item \textbf{Supported platforms:} Any platform with a C99 compiler and
          a cryptographic random number generator.  Optional hardware
          acceleration for AES and SHA is recommended but not required.
\end{itemize}

\section{Testing}

\subsection{Conformance Test Suite}

Implementations MUST pass the SecureLink conformance test suite, which
validates:

\begin{itemize}
    \item Correct encoding and decoding of all message types.
    \item Proper handling of all defined alert codes.
    \item Interoperability with a reference implementation.
    \item Resistance to timing side-channel attacks (constant-time comparison
          for all secret-dependent operations).
\end{itemize}

\subsection{Test Vectors}

The following test vectors are provided for implementation validation:

\begin{verbatim}
ClientNonce:  0x000102030405060708090A0B0C0D0E0F
              101112131415161718191A1B1C1D1E1F
ServerNonce:  202122232425262728292A2B2C2D2E2F
              303132333435363738393A3B3C3D3E3F
CipherSuite:  0x01
ExpectedHash: 4A8F5A2C3B7D9E1F0A2B3C4D5E6F7A8B
              9C0D1E2F3A4B5C6D7E8F9A0B1C2D3E4F
\end{verbatim}

\subsection{Performance Benchmarks}

\begin{table}[htbp]
    \centering
    \caption{Reference performance benchmarks (single core).}
    \label{tab:perf}
    \begin{tabular}{@{}lrrr@{}}
        \toprule
        \textbf{Operation} & \textbf{ARM Cortex-M4} & \textbf{x86-64} &
            \textbf{Target} \\
        \midrule
        Handshake (cold)   & 8.2\,ms   & 1.4\,ms  & $< 10$\,ms \\
        Handshake (resume) & 3.1\,ms   & 0.6\,ms  & $< 5$\,ms  \\
        Encrypt (1\,KB)    & 0.42\,ms  & 0.03\,ms & $< 1$\,ms  \\
        Decrypt (1\,KB)    & 0.44\,ms  & 0.03\,ms & $< 1$\,ms  \\
        \bottomrule
    \end{tabular}
\end{table}

\printbibliography[heading=none]

\end{document}
